%%%%%%%%%%%%%%%%%%%%%%%%%%%%%%%%%%%%%%%
%%   SECTION 5: The Pontryagin $\nu$-invariant of the char map    %%%%%%
%%%%%%%%%%%%%%%%%%%%%%%%%%%%%%%%%%%%%%%


\section{The Pontryagin $\nu$-invariant of the characteristic map} 
\label{sec:Aprep}

Fix  a link $L$ in $S^3$ with three components $X$, $Y$ and $Z$ parametrized by $x=x(s)$, $y=y(t)$ and $z=z(u)$.  Recalling that $S^3$ is regarded as the group of unit quaternions, the asymmetric characteristic map $h_L:\tor\to S^2$ is given by
$$
h_L(s,t,u) \ = \ (y_z-x_z)/|y_z-x_z|
$$
where $a_z$ is an abbreviation for the vector $\pr_{-1}(-a\zbar)$ in pure imaginary quaternion $3\text{-space}$ $\br^3$.   To carry out the proof of \thmref{A}, we need a way to compute the absolute Pontryagin $\nu$-invariant  of $h_L$.  A procedure for doing so is described here.

Throughout this section and the next, $\br^3$ is pictured in the usual way, with the $ij$-plane horizontal and the $k$-axis pointing straight up as in \figref{quat}.  In particular, we view $k$ as the north pole of the unit sphere $S^2$.  

%%% FIGURE 15: ijk-space %%%
\begin{figure}[h!]
\includegraphics[height=150pt]{Figijk.pdf}
\put(-99,40){\small$i$}
\put(-33,58){\small$j$}
\put(-89,117){\small$k$}
\caption{Pure imaginary quaternion $3$-space}
\label{fig:quat}
\end{figure}
%%%%%%%%%%%%%%%%%%%%


%%%%%%%%%%%%%%%%%%%%%%%%%%%
%%%%%%%%%%%%%%%%%%%%%%%%%%%
\part{Outline of the procedure for computing $\nu(h_L)$}
%%%%%%%%%%%%%%%%%%%%%%%%%%%
%%%%%%%%%%%%%%%%%%%%%%%%%%%


First, we will manipulate $L$ by a link homotopy into a favorable position -- referred to as {\itb generic} below -- so that, in particular, the north pole $k$ of $S^2$ is a regular value of $h_L$.  

Then we will determine the associated {\itb Pontryagin link}  
$$
\bl \ = \ h_L^{-1}(k) \ \subset \ \tor
$$
by giving a simple method for constructing a toral diagram for it, in the sense of \secref{pontryagin}.  Roughly speaking, our approach is as follows.  

By definition of $h_L$, the Pontryagin link $\bl$ consists of all $(s,t,u)\in\tor$ for which the vector in $\br^3$ from $x_z$ to $y_z$ points straight up, where $x=x(s)$, $y=y(t)$ and $z=z(u)$.   The genericity of $L$ will imply that for some points $(s,t)$ in the $2$-torus $\torh$ there is a {\it unique} $u = u(s,t)\in S^1$ that will make this happen, while for all other points $(s,t)$, {\it no} $u$ will work.  Furthermore, the set $\cald$ of all points of the first kind, called {\itb isogonal points} for reasons explained below, is a smooth $1$-dimensional  submanifold of $T^2$ whose components we call {\itb icycles}.  

Thus 
$$
\bl \ = \ \{(s,t,u)\in\tor \st (s,t)\in\cald \text{ and } u=u(s,t)\}\,,
$$
the graph of the function $u(s,t)$ over the collection $\cald$ of icycles in $T^2$.   When suitably oriented and decorated with framing and vertical winding numbers, $\cald$ will be the desired toral diagram of $\bl$.  


In particular, the icycles in $\cald$ correspond to certain oriented cycles of vectors directed from $X$ to $Y$, which we call {\itb bicycles}, that are easily spotted from a picture of $L$.  Each bicycle has a {\itb longitudinal} and {\itb meridional degree}, recording how much it turns and spins relative to the standard open book structure on $S^3$.  The framing and vertical winding number of the corresponding icycle are determined by these degrees. 

Therefore a diagram for $\bl$ can be constructed once we identify the bicycles in $L$.  With this diagram in hand, the methods of \secref{pontryagin} can then be used to compute $\nu(h_L)$.

We now give the details of this procedure.  There are three 
geometries involved: spherical geometry in $S^3$,  where the link  $L$ lives, Euclidean geometry in $\br^3$,  the setting for the asymmetric characteristic map, and hyperbolic geometry in the complex upper half plane $\bh$,  which turns out to be for us the natural geometry on the pages of the ``standard" open book decomposition of $\br^3$.  We begin with an explicit construction of this open book, which provides a framework for the discussion that follows.  

\begin{comment}
Our aim is first to give a working description of the procedure, then to give some examples, and to end with the proofs.
\end{comment}

%\break

%%%%%%%%%%%%%%%%%%%%%%%%%%%
%%%%%%%%%%%%%%%%%%%%%%%%%%%
\part{The standard open books in $\br^3$ and $S^3$}
%%%%%%%%%%%%%%%%%%%%%%%%%%%
%%%%%%%%%%%%%%%%%%%%%%%%%%%

Consider the great circle $K$ in $S^3$ through $1$ and $k$, and the orthogonal great circle $C$ through $i$ and $j$.  Orient both circles by left complex multiplication by $K$ (i.e.\ from $1$ toward $k$ on $K$, and $i$ toward $j$ on $C$) so that their linking number is $+1$.

Stereographic projection from $-1$ carries $K$ 
% (punctured at $-1$) 
onto the $k$-axis in $\br^3$, and fixes $C$, which now appears as the unit circle in the $ij$-plane.  The complement $\V$ of the $k$-axis in $\br^3$ is naturally identified with the product of a circle $S^1$ (viewed as the quotient $\br/2\pi\bz$) with the complex upper half plane $\bh$ (which for later purposes will be viewed as the hyperbolic plane) via the diffeomorphism
% \V =  \br^3-k\text{-axis}
$$
\V \ \longrightarrow \ S^1 \times \bh
$$ 
with coordinates \,$\ell: \V\longrightarrow S^1$\, and \,$m: \V\longrightarrow\bh$\, given by 
$$
\ell(q) \ = \ \arg(q_1+q_2i)\qquad\text{and}\qquad m(q) \ = \ q_3 + |q_1+q_2i|i\,.
$$
for $q=q_1i+q_2j+q_3k\in \V =  \br^3-k\text{-axis}$.  In other words, if $q=(r,\theta,z)$ in cylindrical coordinates, then $\ell(q) = \theta$ and $m(q) = z+ri$.  We call $\ell$ and $m$ the {\itb longitudinal} and {\itb meridional projections} in $\br^3$, and refer to $\ell(q)$ as the {\itb polar angle} of $q$.  
% Two points with the same polar angle will be called {\itb isogonal}.  


The longitudinal projection defines the {\itb standard open book} in $\br^3$, with binding the $k$-axis, and with pages $P_\theta = \ell^{-1}(\theta)$ for $\theta\in S^1$.  The pages are just the oriented vertical half-planes bounded by the $k$-axis, each indexed by its polar angle as shown in \figref{book}.  The meridional projection serves to identify each page $P_\theta$ with the hyperbolic plane $\bh$, with ``center" \,$i_\theta = i\,\cos\theta + j\,\sin\theta$\, corresponding to $i$.  The union of all the page centers is the unit circle $C$.

%%% FIGURE 16: ijk-space %%%
\begin{figure}[h!]
\includegraphics[height=170pt]{FigBook.pdf}
\put(-113,63){\small$i$}
\put(-64,66){\small$i_\theta$}
\put(-49,82){\small$j$}
\put(-89,114){\small$k$}
\put(-90,43){\small$\theta$}
\put(-107,-10){binding}
\put(-145,142){pages}
\put(-53,10){$P_\theta$}
\put(-133,62){$C$}
\caption{The standard open book in $\br^3$}
\label{fig:book}
\end{figure}
%%%%%%%%%%%%%%%%%%%%

\break

The longitudinal and meridional projections  in $\br^3$ lift to $S^3$ \label{longproj} (in the complement of $K$) by composing with stereographic projection from $-1$.  Relying on the context, we continue to denote them by $\ell$ and $m$, and to refer to $\ell(q)$ as the polar angle of $q$.  Explicitly
$$
\ell(q) \ = \ \arg(q_1+q_2i)\qquad\text{and}\qquad m(q) \ = \ \frac{q_3 + |q_1+q_2i|i}{1+q_0}\,.
$$
for $q = q_0+q_1i+q_2j+q_3k\in S^3-K$.  Here we are using 
the formula 
$$
\pr_{-1}q \ =  \ (\Im q) / (1 + \re q)  =  (q_1 i + q_2 j + q_3 k) / (1 + q_0)\,,
$$
and then plugging this into the formula above for the longitudinal and meridional projections in $\br^3$.  

Just as in $\br^3$, the longitudinal projection in $S^3$ defines the {\itb standard open book} in $S^3$, with binding $K$, and with pages $H_\theta = \ell^{-1}(\theta)$.  The pages are now open great hemispheres  in $S^3$ with centers (i.e.\ poles) $i_\theta$ along the great circle $C$, and with $K$ as equator.  By design, $\pr_{-1}$ maps each hemispherical page $H_\theta$ in $S^3$ onto the corresponding half-planar page $P_\theta$ in $\br^3$.   


We now turn our attention back to links in the $3$-sphere. 


%%%%%%%%%%%%%%%%%%%%%%%%%%%
%%%%%%%%%%%%%%%%%%%%%%%%%%%
\part{Generic links}
%%%%%%%%%%%%%%%%%%%%%%%%%%%
%%%%%%%%%%%%%%%%%%%%%%%%%%%

A three-component link $L=X\,\cup\,Y\,\cup\,Z$ in $S^3$ is {\itb generic} if

\vspace{.02in}
\begin{indentation}{.05em}{1em}
\begin{enumerate}
\item $Z$ coincides with the oriented binding $K$ of the standard open book, and
\smallskip
\item $X$ and $Y$ wind ``generically" around $Z$.
\end{enumerate}  
\end{indentation}
\vspace{.02in}

More precisely, the second condition requires the restriction to $X\,\cup\,Y$ of the longitudinal projection \,$\ell:S^3-K\longrightarrow S^1$, which sends each point to its polar angle, to be a Morse function with just one critical point per critical value.  Geometrically, this means $X$ and $Y$ are transverse to the pages of the standard open book in $S^3$, except for finitely many pages where exactly one of them turns around at a single point.   These will be called the {\itb critical points} of $L$, while all other points on $X\cup Y$ will be called {\itb regular points}.   Each regular point $w$ has a sign, denoted {\bf\boldmath $\sgn(w)$}, when viewed as an intersection point of $L$ with the page containing $w$.  Thus $\sgn(w) = +1$ or $-1$ according to whether $L$ is oriented in the direction of increasing or decreasing polar angle near $w$.  

An example of a generic link is shown in \figref{generic}.  It has four critical points, two on $X$ and two on $Y$, indicated by dots in the picture.


%%% FIGURE 17: Generic Link %%%
\begin{figure}[h!]
\includegraphics[height=150pt]{FigGeneric.pdf}
\put(-150,110){$X$}
\put(-152,30){$Y$}
\put(-90,138){$Z$}
\caption{A generic link}
\label{fig:generic}
\end{figure}
%%%%%%%%%%%%%%%%%%%%

Any three-component link in $S^3$ is evidently link homotopic to a generic one:  first unknot the last component by a link homotopy and move it to coincide with $K$, and then adjust the first two components by a small isotopy to satisfy the genericity condition (2). 

\smallskip

{\it From this point on we assume that $L$ is generic without further mention.}



\break

As a consequence, we will show the following:

\vspace{.02in}
\begin{indentation}{.05em}{1em}
\begin{enumerate}
\item[(a)] The north pole $k\in S^2$ is a regular value of the characteristic map $h_L$.
\smallskip
\item[(b)] There is a simple method for constructing a toral diagram for the associated Pontryagin link $\bl = h_L^{-1}(k)$ from a picture of $L$.
\end{enumerate}  
\end{indentation}
\vspace{.02in} 

This is the content of the ``bicycle theorem" below.  To state it precisely, we need to introduce the key notion of a {\it bicycle} in $L$, and its associated {\it icycle} in $T^2$.


%%%%%%%%%%%%%%%%%%%%%%%%%%%
%%%%%%%%%%%%%%%%%%%%%%%%%%%
\part{Bicycles and icycles}
%%%%%%%%%%%%%%%%%%%%%%%%%%%
%%%%%%%%%%%%%%%%%%%%%%%%%%%

Assume, as always, that the components $X$, $Y$ and $Z$ of $L$ are parametrized by smooth functions $x=x(s)$, $y=y(t)$ and $z=z(u)$ with nowhere vanishing derivatives.   In particular, points $(s,t)$ in the $2$-torus parametrize pairs of points $x\in X$ and $y\in Y$.  Suppose that $x$ and $y$ have the same polar angle $\theta$, or equivalently that they lie in a common page $H_\theta$ of the standard open book in $S^3$.  Then we call $(s,t)$ an {\itb isogonal point} in $T^2$, and call $(x,y)$ a {\itb page vector} in $L$ with polar angle $\theta$, reflecting the fact that the vector in $\br^3$ from $\pr_{-1}(x)$ to $\pr_{-1}(y)$ lies entirely on the half-planar page $P_\theta$.  

A page vector $(x,y)$ will be called {\itb critical} if $x$ or $y$ is a critical point of $L$, and {\itb regular} if both $x$ and $y$ are regular.  A regular page vector is {\itb positive} if the oriented strands of $L$ through $x$ and $y$ point in the same direction, meaning $\sgn(x) = \sgn(y)$, and is {\itb negative} if they point in opposite directions.  These notions are illustrated in \figref{pagevectors}, in which the vectors labeled $1$ and $2$ are positive regular page vectors, $3$ is negative regular, and $4$ is critical.

%%% FIGURE 18: Page vectors %%%
\begin{figure}[h!]
\includegraphics[height=170pt]{FigPageVectors2.pdf}
%\put(-94,67){\small$\vec_1$}
%\put(-93,30){\small$\vec_2$}
%\put(-65,45){\small$\vec_3$}
%\put(-40,60){\small$\vec_4$}
\put(-150,105){\tiny$1$}
\put(-156,73){\tiny$2$}
\put(-112,85){\tiny$3$}
\put(-65,98){\tiny$4$}
\put(-185,105){$X$}
\put(-20,102){$Y$}
\put(-106,18){$Z$}
\caption{Page vectors}
\label{fig:pagevectors}
\end{figure}
%%%%%%%%%%%%%%%%%%%%

Now consider the spaces
$$
\cald \ = \ \{\text{isogonal points in $T^2$}\} \qquad\text{and}\qquad \calp \ = \ \{\text{page vectors in $L$}\} \,.
$$
By definition $\cald$ parametrizes $\calp$.  The genericity of $L$ implies that $\calp$ consists of a finite number of disjoint cycles of page vectors, and that $\cald$ consists of a finite collection of smooth simple closed curves in $T^2$.   Orient $\calp$ so that it points to the right (meaning in the direction of increasing polar angle) at each positive regular page vector in it, and to the left at each negative regular page vector.  This gives a well-defined orientation on $\calp$, inducing one on $\cald$ as well.  We call these the {\itb preferred orientations} on $\calp$ and $\cald$.     

\begin{definition}
A {\itb bicycle} (or ``bi-cycle")  in $L$ is a component $\calp_i$ of $\calp$, that is, an oriented cycle of page vectors.  Each bicycle is parametrized by a component $\cald_i$ of $\cald$, which we call its associated {\itb icycle} (or ``i-cycle").
\end{definition}


%%%%%%%%%%%%%%%%%%%%%%%%%%%
%%%%%%%%%%%%%%%%%%%%%%%%%%%
\part{Some examples of bicycles and their associated icycles}
%%%%%%%%%%%%%%%%%%%%%%%%%%%
%%%%%%%%%%%%%%%%%%%%%%%%%%%

We first draw in \figref{local} four local pictures of a bicycle near a regular page vector, and below them, their parametrizing icycles.  It is understood that these pictures take place somewhere in front of the upward pointing $Z$ axis. The four cases represent the possible directions of  $X$ and $Y$ relative to the page containing the vector.  In each case, the orientation of the bicycle is indicated by a squiggly arrow.
%%% FIGURE 19: Local Bicycle %%%
\begin{figure}[h!]
\includegraphics[height=175pt]{FigLocal.pdf}
\put(-374,146){\small$Y$}
\put(-376,96){\small$X$}
%\put(-347,147){\small$y$}
%\put(-348,98){\small$x$}
\put(-215,147){\small$Y$}
\put(-217,96){\small$X$}
%\put(-245,148){\small$y$}
%\put(-245,99){\small$x$}
\put(-115,147){\small$Y$}
\put(-171,96){\small$X$}
%\put(-143,148){\small$y$}
%\put(-142,99){\small$x$}
\put(-65,147){\small$Y$}
\put(-13,96){\small$X$}
%\put(-38,148){\small$y$}
%\put(-39,99){\small$x$}
\put(-318,3){\scriptsize$s$}
\put(-363,50){\scriptsize $t$}
\put(-216,3){\scriptsize$s$}
\put(-261,50){\scriptsize $t$}
\put(-114,2){\scriptsize$s$}
\put(-159,49){\scriptsize $t$}
\put(-12,2){\scriptsize$s$}
\put(-57,49){\scriptsize $t$}
\caption{Local pictures of a bicycle and its corresponding icycle}
\label{fig:local}
\end{figure}
%%%%%%%%%%%%%%%%%%%% 


For an example of a full bicycle, consider the ``clasp" between $X$ and $Y$ pictured in \figref{claspbike}(a).  This gives rise to the bicycle in \figref{claspbike}(b), passing successively through the vectors labeled $1,2,3,4$ and then back to $1$.  The associated icycle is a counterclockwise circle in $T^2$, shown in \figref{claspbike}(c).  The route taken by this bicycle is ``short" in the sense that it does not wind around the binding, although it does spin within the pages.

%%% FIGURE 20: Clasp/Bicycle %%%
\begin{figure}[h!]
\includegraphics[height=80pt]{FigClasp.pdf}
\put(-280,53){\small$X$}
\put(-350,15){\small$Y$}
\put(-325,66){\small$Z$}
\put(-198,45){\tiny$1$}
\put(-160,51){\tiny$2$}
\put(-202,33){\tiny$3$}
\put(-160,25){\tiny$4$}
\put(2,6){\small$s$}
\put(-82,73){\small$t$}
\put(-42,64){\tiny$1$}
\put(-65,40){\tiny$2$}
\put(-42,17){\tiny$3$}
\put(-19,40){\tiny$4$}
\put(-340,-20){(a) the clasp}
\put(-215,-20){(b) the bicycle}
\put(-80,-20){(c) the icycle}
\caption{Bicycle arising from a clasp}
\label{fig:claspbike}
\end{figure}
%%%%%%%%%%%%%%%%%%%% 

As another example, the link shown in \figref{generic} and reproduced in \figref{icycles}(a) below has three bicycles.  Two of them are short, arising from the clasps as in the previous example, while the remaining long one oscillates back and forth in the longitudinal direction, eventually making one full revolution around the binding.  It is an instructive exercise left to the reader to recover the plot of the associated icycles in \figref{icycles}(b), in which the trivial circles labeled $A$ and $B$ correspond to the clasps in $L$ with the same labels, $C$ labels the icycle that parametrizes the long bicycle, and the corners of the square parametrize the pair $(x,y)$ indicated by the dots in \figref{icycles}(a).   

%%% FIGURE 21: Icycles %%%
\begin{figure}[h!]
\includegraphics[height=125pt]{FigIcycles.pdf}
\put(2,20){\small$s$}
\put(-83,104){\small$t$}
\put(-300,90){\small$X$}
\put(-295,20){\small$Y$}
\put(-238,115){\small$Z$}
\put(-248,54){\tiny$A$}
\put(-215,52){\tiny$B$}
\put(-44,53){\tiny$A$}
\put(-29,36){\tiny$B$}
\put(-27,80){\tiny$C$}
\put(-280,-15){(a) the generic link $L$}
\put(-90,-15){(b) the icycles of $L$}
\caption{A generic link and its icycles}
\label{fig:icycles}
\end{figure}
%%%%%%%%%%%%%%%%%%%%

Before proceeding, we remind the reader that our interest in icycles associated to $L$ stems from the fact that -- when suitably decorated -- they give a toral diagram for a Pontryagin link of the characteristic map $h_L$.   This is the content of the bicycle theorem.  
% To state it, we must first identify the characteristics of bicycles that determine these decorations.

%%%%%%%%%%%%%%%%%%%%%%%%%%%
%%%%%%%%%%%%%%%%%%%%%%%%%%%
\part{The Bicycle Theorem}
%%%%%%%%%%%%%%%%%%%%%%%%%%%
%%%%%%%%%%%%%%%%%%%%%%%%%%%

The longitudinal and meridional projections on $S^3-K$, defined earlier,
% on page \pageref{longproj}
induce projections by the same name on the space $\calp$ of page vectors,
$$
S^1 \, \overset{\ell}{\longleftarrow} \, \calp \, \overset{m}{\longrightarrow} \, \bh \,,
$$
given by $\ell(x,y) = \ell(x) = \ell(y)$ and $m(x,y) = \arg(m(y)-m(x))$.  In other words $\ell(x,y)$ is the polar angle that parametrizes the common hemispherical page in $S^3$ containing $x$ and $y$, and $m(x,y)$ is the argument of the vector from $m(x)$ to $m(y)$ in 
% the upper-half plane 
$\bh$.  

Using these projections, we define the {\itb longitudinal} and {\itb meridional degrees} of a bicycle $\calp_i$ in $L$ by  
$$
\ell_i \ = \ \deg(\ell|\calp_i) \qquad\text{and}\qquad  m_i \ = \ \deg(m|\calp_i). \label{degrees}
$$
These integer invariants record, respectively, the number of times $\calp_i$ travels around the binding $Z$, and the number of times its vectors spin around in the pages as it goes.   

For example, any bicycle arising from a clasp between $X$ and $Y$ has zero longitudinal degree, while its meridional degree can be $\pm1$.  In particular, the one shown in \figref{claspbike} has meridional degree $-1$, while the ones labeled $A$ and $B$ in \figref{icycles}(a) have degrees $1$ and $-1$, respectively.  The long bicycle in \figref{icycles}(a), labeled $C$, has longitudinal degree $2$ and meridional degree $1$.  

For any icycle $\cald_i$ in $T^2$, parametrizing a bicycle $\calp_i$ in $L$, define the {\itb framing} $n_i$ and {\itb vertical winding number} $r_i$ of $\cald_i$ by
$$
n_i \ = \ -\ell_i-m_i \qquad\text{and}\qquad r_i \ = \ m_i 
$$
where $\ell_i$ and $m_i$ are the longitudinal and meridional degrees of $\calp_i$.

We can now state the main result of this section.

%%%%% BICYCLE THEOREM  %%%%%%
\begin{theorem*}\label{thm:bicycle}
\textbf{\boldmath Let $L$ be a generic link in $S^3$.  Then
\begin{indentation}{-.2em}{1.2em}
\begin{itemize}
\item[(a)] The north pole $k\in S^2$ is a regular value of $h_L:\tor\to S^2$.
\smallskip  
\item[(b)] The collection $\cald$ of icycles of $L$, together with their framings and vertical winding numbers as defined above, forms a toral diagram for the associated Pontryagin link $\bl = h_L^{-1}(k)$.
\end{itemize}
\end{indentation}}
\end{theorem*}
%%%%%%%%%%%%%%%%%%


Before proving this theorem, we illustrate how it can be used to compute the Pontryagin invariant of the characteristic map of a generic link.

%%%%%%%%%%%%%%%%%%%%%%%%%%%
%%%%%%%%%%%%%%%%%%%%%%%%%%%
\part{Computing $\nu(h_L)$ for a generic link $L$ using the bicycle theorem}
%%%%%%%%%%%%%%%%%%%%%%%%%%%
%%%%%%%%%%%%%%%%%%%%%%%%%%%



%%%%%%%%%%%%%%%%%%%%%%%%%%%%%%%%%%%%%%%%%%%%%%%%
%%%%%%%%%%%%%%%%%%%%%%%%%%%%%%%%%%%%%%%%%%%%%%%%
%%%%%%%%%%%%%%%%%%%%%%%%%%%%%%%%%%%%%%%%%%%%%%%%
\begin{comment}

Since the diagram in the lemma has no isolated marked points, and the internal marked points are recorded by the vertical winding numbers of the icycles, $\nu(h_L)$ can be computed using \propref{pontryagin} once the depths of the icycles are determined.  As explained in \S3, this requires the choice of a base point $(s_0,t_0)$ in $\torh$ lying in the complement of $\cald$, base points $(s_i,t_i)$ on each icycle $\cald_i$,  and paths $\gamma_i$ from points $(s_i',t_i')$ ``adjacent" to $(s_i,t_i)$ (displaced slightly to the right of $\cald_i$) to $(s_0,t_0)$.  Then the depth $d_i$ of $\cald_i$ is equal to the intersection number $ \gamma_i\cdot\cald$, which is well-defined modulo the greatest common divisor of the horizontal winding numbers of the diagram, and easily computed from a sketch of $\cald$.

Alternatively, $d_i$ can be computed directly from a picture of $L$ without having to sketch $\cald$, as follows.  The base point in $T^2$ corresponds to a pair of base points $x_0\in X$ and $y_0\in Y$ lying on {\it distinct} hemispherical pages $H_0$ and $K_0$ in $S^3$, and the base points on the icycles correspond to a choice of page vectors $(x_i,y_i)$ in the corresponding bicycles.  We can assume that these page vectors are all regular, and that the hemispherical pages $H_i$ containing $x_i$ and $y_i$ do not contain either base point $x_0$ or $y_0$.  Then $X-\{x_0,x_i\}$ consists of two open oriented arcs; let $\alpha_i$ denote either one of these arcs.  Similarly, let $\beta_i$ denote either one of the two arcs in $Y-\{y_0,y_i\}$.  Then a simple geometric argument shows that 
$$
d_i = \alpha_i\cdot H_i - \beta\cdot H_0 - \textstyle\frac12(\sgn(x_i)-1) \mod{\gcd(p,q)}
$$
where $\sgn(x_i)$ denotes the intersection number with $H_i$ of a short arc in $X$ containing $x_i$.  This formula is derived by choosing $\gamma_i$ to be the ``rectilinear" L-shaped path $\gamma_i = [s_i',s_0]\times t_i' \,\cup\, s_0\times [t_i',t_0]$ in $T^2$. 

\end{comment}
%%%%%%%%%%%%%%%%%%%%%%%%%%%%%%%%%%%%%%%%%%%%%%%%
%%%%%%%%%%%%%%%%%%%%%%%%%%%%%%%%%%%%%%%%%%%%%%%%
%%%%%%%%%%%%%%%%%%%%%%%%%%%%%%%%%%%%%%%%%%%%%%%%


As a first example, again consider the link $L$ pictured in \figref{icycles}(a).   As noted above, it has three bicyles $A$, $B$, $C$, with longitudinal degrees $0,0,2$, meridional degrees $1,-1,1$, and so by definition, framings $-1,1,-3$ and vertical winding numbers $1,-1,1$.   

By the bicycle theorem, the Pontryagin link for the characteristic map $h_L$ has toral diagram as shown in \figref{icycles}(b) with vertical winding numbers $1$, $-1$ and $1$ on the icycles $A$, $B$ and $C$, and with global framing $n=-1+1-3=-3$.  Thus the total vertical winding number is $r = 1-1+1=1$ and from the diagram we compute the horizontal winding numbers to be $p=-1$ and $q=-2$.  (These values for the winding numbers of the diagram are confirmed by the calculations $p = \lk(Y,Z)=-1$, $q = \lk(X,Z)=-2$ and $r = \lk(X,Y)=1$.)  Thus the invariant $\nu(h_L)$ is well defined modulo $2 = 2\gcd(-1,-2,1)$.

Using a base point in the lower right corner of the diagram, and straight line paths from the icycles to the base point, the depths of the icycles $A$, $B$ and $C$ are $1$, $-1$ and $-1$.  Thus by \propref{pontryagin} we conclude that 
$$
\nu(h_L) \ = \ n+pq + \textstyle\sum d_ir_i \ = \ -3+2+1  \ = \ 0 \ \in \ \bz_2.
$$

Although the purpose of this example is to illustrate how the bicycle theorem is used for computations, we note that \thmref{A} (yet to be proved) yields the same result here effortlessly, since it implies that 
the Pontryagin invariant of the characteristic map of any three-component link in $S^3$ is even.  Therefore, when the pairwise linking numbers $p$, $q$ and $r$ are relatively prime, as they are in this case, the computation is  modulo $2\gcd(p, q, r)  =  2$,  and so the Pontryagin invariant is zero.   

%%%%%%%%%%%%%%%%%%%%%%%%%%%
%%%%%%%%%%%%%%%%%%%%%%%%%%%
\part{Double crossing changes}
%%%%%%%%%%%%%%%%%%%%%%%%%%%
%%%%%%%%%%%%%%%%%%%%%%%%%%%

For our next example we analyze the effect on $\nu(h_L)$ of changing two crossings of opposite signs between the first two components of a generic link $L=X\cup Y\cup Z$.  This will be a key step in our inductive proof of \thmref{A}.

To this end, choose a positive and a negative crossing between $X$ and $Y$ in a suitable projection of $L$, and let $(x_+,y_+)$ and $(x_-,y_-)$ be the corresponding page vectors.  Changing both of the crossings yields a new link $\what L$, with the same pairwise linking numbers as $L$.  This is illustrated in \figref{dblbor} where $L$ is the Borromean rings, shown on the left, and $\what L$ is the unlink, shown on the right.  

%%% FIGURE 22: Double crossing %%%
\begin{figure}[h!]
\includegraphics[height=135pt]{FigDblBor.pdf}
\put(-348,97){$X$}
\put(-335,17){$Y$}
\put(-325,125){$Z$}
\put(-360,-15){(a) the Borromean rings $L$}
\put(-45,97){$\what X$}
\put(-60,17){$\what Y$}
\put(-88,125){$\what Z$}
\put(-115,-15){(b) the unlink $\what L$}
\caption{A double crossing change}
\label{fig:dblbor}
\end{figure}
%%%%%%%%%%%%%%%%%%%%

We then say that $\what L$ is obtained from $L$ by a {\itb double crossing change}, and propose to use the bicycle theorem to compute the resulting change
$$
\Delta\nu \ = \  \nu(h_{\what L}) - \nu(h_{L})
$$
in Pontryagin invariants.  


As it turns out, there is a simple formula for $\Delta\nu$ involving the link $L_0$ obtained from $L$ by ``smoothing" both crossings in the usual way: 
\vskip .1pt
$$
\includegraphics[height=20pt]{FigSmoothCrossing.pdf}
$$
\vskip .1pt
The smoothed link $L_0$ has three components, $P$ and $Q$ (replacing $X$ and $Y$) and $Z$.  The component $P$ goes from $+$ to $-$ along $Y$, and then back from $-$ to $+$ along $X$, where we retain the $+$ and $-$ labels after smoothing, while $Q$ goes from $+$ to $-$ along $X$, and then back from $-$ to $+$ along $Y$.  This is illustrated in \figref{smooth} for the double crossing change shown in \figref{dblbor}, where we use  $X_{+-}$ to denote the arc on $X$ from $+$ to $-$\,, and so forth.  

%%% FIGURE 23: Smoothing a double crossing %%%
\begin{figure}[h!]
\includegraphics[height=135pt]{FigSmooth.pdf}
\put(-23,60){$P$}
\put(0,50){$Q$}
\put(-90,125){$Z$}
\put(-103,99){\footnotesize$X_{+-}$}
\put(-65,74){\footnotesize$X_{-+}$}
\put(-95,61){\footnotesize$Y_{+-}$}
\put(-67,21){\footnotesize$Y_{-+}$}
%%%%%%%%%%%
\caption{The smoothed link $L_0 \ = \ P\cup Q\cup Z$}
\label{fig:smooth}
\end{figure}
%%%%%%%%%%%%%%%%%%%%

Then we have the following consequence of the Bicycle Theorem:


\begin{corollary}\label{cor:double}
\textbf{\boldmath \textup{(Double Crossing)}
If $L = X\cup Y\cup Z$ is transformed into $\what L$ by a double crossing change, and $P$ and $Q$ are the components of the associated smoothing of $X\cup Y$, as explained above, then the corresponding Pontryagin invariants change by 
$$
\Delta\nu \ = \ 2\,\lk(P,Z) \ = \ -2\,\lk(Q,Z) \ \in \ \bz_{2\gcd(p,q,r)}
$$
where $p$, $q$ and $r$ are the pairwise linking numbers of the components of $L$.}
\end{corollary}


\begin{proof}
First note that
$$
\lk(P\cup Q,Z) \ = \  \lk(X\cup Y,Z) \ = \ p+q \ = \ 0 \ \in \ \bz_{\gcd(p,q,r)}
$$
so it suffices to establish the first equality in the corollary.

We may assume that the two page vectors  $\vv_+ = (x_+,y_+)$ and $\vv_- = (x_-,y_-)$ associated with the chosen crossings lie on distinct hemispherical pages $H_+$ and $H_-$ of the standard open book in $S^3$, and that these pages contain no critical points in the link.   These two page vectors may lie in distinct bicycles in $L$, or they may lie in the same bicycle.

Suppose first that $\vv_+$ and $\vv_-$  lie in distinct bicycles $\calp_+$ and $\calp_-$ in $L$.  When we change $L$ to $\what L$, these bicycles will change, but the icycles  $\cald_+$ and $\cald_-$ in the $2$-torus that parametrize them will stay the same.  However,  their vertical winding numbers  $r_+$ and $r_-$ (which record the meridional degrees of $\calp_+$ and $\calp_-$) and their framings $n_+$ and $n_-$ (which record the negative of the sum of the meridional and longitudinal degrees of $\calp_+$ and $\calp_-$) will change.  In particular, we claim that 
$$
\what r_+\ = \ r_ + -1 \qquad\text{and}\qquad \what r_- \ = \ r_- + 1\,,
$$
and consequently \, $\what n_+ = n_+ + 1$ \,and\, $\what n_- = n_- - 1$ \, since the longitudinal degrees of the bicycles clearly do not change.  The figure below helps us to see this.
$$
\includegraphics[height=30pt]{FigVert.pdf}
\put(-212,0){\small$X$}
\put(-209,25){\small$Y$}
\put(-175,-1){\small$L$}
\put(-77,0){\small$\what X$}
\put(-75,25){\small$\what Y$}
\put(-40,-2){\small$\what L$}
$$
\vskip .2pt
In this figure, we start with a positive crossing in $L$ and change it to a negative crossing in $\what L$. Since the strands of $X$ and $Y$ are both pointing to the right, the bicycle $\calp_+$ is moving from left to right.  During this motion the page vectors in $\calp_+$  undergo half a {\it counter-clockwise} rotation with respect to the preferred orientation on the pages.  In the corresponding picture for $\what L$ , we see half a {\it clockwise} rotation.  Therefore the bicycle in $\what L$ has one more full clockwise rotation in the pages than in $L$, and so $\what r_+ =  r_+ -1$ as claimed.

If, for example, we switched the arrow on the strand of $Y$, we would have a negative crossing, but then the corresponding bicycle would be moving from right to left, and we would see that $\what r_- = r_- +1$ as claimed.  With this guidance, we leave the remaining cases to the reader.

Now suppose that the page vectors $\vv_+$ and $\vv_-$ lie in the same bicycle $\calp_+ = \calp_-$.  Then the above changes in vertical winding number $r$ for the corresponding icycle $\cald_+=\cald_-$ will cancel, and so we see that neither the vertical winding number nor the framing of the icycle will change, that is $\what r_\pm = r_\pm$ and $\what n_\pm = n_\pm$.

Note that in either case, whether $\vv_+$ and $\vv_-$ lie in the same or different bicycles, the {\it total framing} $n$ of the diagram (the sum of the framings of the icycles) does not change.

At this point we recall \propref{pontryagin}, which tells us that 
$$ 
\nu(h_L)  \ = \  n  +  p q  +  \textstyle\sum d_i r_i\,. 
$$
In passing from $L$  to  $\what L$ via the double crossing change, we have just seen that the total framing $n$ does not change, and the pairwise linking numbers $p$ and $q$ certainly do not change.  We have also seen that the icycles stay the same, so their depths $d_i$ do not change.  Only the winding numbers $r_+$ and $r_-$ of the (possibly equal) icycles $\cald_+$ and $\cald_-$ may change.  In particular, they also do not change when the page vectors
$$
\vv_+ \ = \ (x(s_+), y(t_+)) \qquad\text{and}\qquad \vv_- \ = \ (x(s_-), y(t_-))
$$ 
lie in the same bicycle, and so $\Delta\nu = 0$ in this case, while they change to $\what r_+ = r_+ - 1$ and $\what r_- = r_- + 1$ when $\vv_+$ and $\vv_-$ lie in distinct bicycles, in which case we have
$$
\begin{aligned}
\Delta\nu \ &= \  \nu(h_{\what L}) - \nu(h_L) \\
\ &= \ d_-(\what r_- - r_-) + d_+(\what r_+ - r_+) \\
\ &= \ d_- - d_+\,.
\end{aligned}
$$
Thus in either case it remains to prove that $d_- - d_+ = 2\,\lk(P,Z)$. 

To show this, we must compute the depths $d_+$ and $d_-$ of $\cald_+$ and $\cald_-$.  This requires a choice of base point, and then a choice of paths $\gamma_+$ and  $\gamma_-$ from $\cald_+$ and $\cald_-$ to this base point.  We let $(s_+,t_-)$ be the base point, and use the vertical path $\gamma_+$ from $(s_+,t_+)$ to $(s_+,t_-)$ and the horizontal path $\gamma_-$ from $(s_-,t_-)$ to $(s_+,t_-)$.  

The depth $d_-$ of the component $\cald_-$ of $\cald$ counts the intersections of $\gamma_-$ with $\cald$, which means that it counts the times that $(x(s),y(t_-))$ is a page vector for $s_-\le s\le s_+$, assuming parametrizations set up so that $s_-<s_+$\,.  Since $y(t_-)$ already lies in the page $H_-$, this means we are counting the times that $x(s)$ also lies in $H_-$ for $s_-\le s\le s_+$\,.  In other words we are counting the intersection number $X_{-+}\dt H_-$\,.  Examining \figref{local}, we find that the signs of the intersection  points of $\gamma_-$ with $\cald$ {\it agree} with the signs of the corresponding intersection points of $X_{-+}$ with $H_-$, and hence
$$
d_- \ = \ 2\, X_{-+} \dth H_- \mod{2\gcd(p,q,r)}.
$$


In a similar fashion, we find that intersections of $\gamma_+$ with $\cald$ correspond to intersections of $Y_{+-}$ with $H_+$, but in this case, the signs of corresponding points of intersection are {\it opposites}, and so the depth of $\cald_+$ is 
$$
d_+ \ = \ -2\,Y_{+-} \dth H_+\mod{2\gcd(p,q,r)}.
$$ 

Therefore the assertion that  $d_- - d_+ = 2\lk(P,Z)$, which will complete the proof of the corollary, reduces to the identity
$$
X_{-+} \dth H_-  +  Y_{+-} \dth H_+ \ = \ \lk(P,Z)\,.
$$
But this follows easily from the fact that $P = X_{-+} \cup Y_{+-}$.  \figref{linking} helps us to see this.  In the figure, the arc  $X_{-+}$  winds around the vertical $Z$ axis $k$ times and the arc $Y_{+-}$ winds around it $\ell$ times, while the entire loop $P = X_{-+} \cup Y_{+-}$ winds around it $k+\ell+1$ times.   By our half-counting convention, we have 
$$
X_{-+} \dt H_- \ = \ 1/2 + k \qquad\text{and}\qquad Y_{+-} \dt H_+ \ = \ 1/2 + \ell
$$
and the result follows.  
\end{proof} 


%%% FIGURE 22: Linking %%%
\begin{figure}[h!]
\includegraphics[height=190pt]{FigLinking.pdf}
\put(-132,10){$H_-$}
\put(-18,12){$H_+$}
\put(-57,27){$X_{-+}$}
\put(-96,155){$Y_{+-}$}
\put(-82,180){$Z$}
\put(-126,50){$\scriptstyle k$}
\put(-17,135){$\scriptstyle \ell$}
\put(-121,86){\boldmath$\scriptstyle -$}
\put(-19,90){\boldmath$\scriptstyle +$}
\caption{The linking of $P=X_{-+} \cup Y_{+-}$ and $Z$}
\label{fig:linking}
\end{figure}
%%%%%%%%%%%%%%%%%%%%


For the inductive step of the proof of \thmref{A} in the next section, we will need to apply this formula for $\Delta\nu$ under the double crossing change $L\to\what L$ shown in \figref{double} (which will be seen to be equivalent to a delta move).  It is understood that $L$ and $\what L$ should coincide outside the picture, where in fact they can be arbitrary.  Indeed, if not generic outside the ball, they can be adjusted by a link homotopy to become so, and then the methods described above apply.   Since the component $Q$ of the smoothed link $L_0$ is just a meridian of $Z$ with $\lk(Q,Z) = -1$, it follows that $\Delta\nu = 2$.  

Thus we have proved the following: \label{double}

%%% FIGURE 25: Double %%%
\begin{figure}[h!]
\includegraphics[height=100pt]{FigDouble.pdf}
\put(-315,14){\small$X$}
\put(-203,73){\small$Y$}
\put(-267,90){\small$Z$}
\put(-300,37){\footnotesize$+$}
\put(-218,37){\footnotesize$-$}
\put(-100,37){\footnotesize$-$}
\put(-18,37){\footnotesize$+$}
\put(-285,-15){(a) the link $L$}
\put(-90,-15){(b) the link $\what L$}
\caption{A simple double crossing change}
\label{fig:double}
\end{figure}
%%%%%%%%%%%%%%%%%%%%

\begin{corollary}\label{cor:double}
\textbf{\boldmath 
If two links $L$ and $\what L$ coincide outside a $3$-ball, and appear in the ball as shown in \figref{double},  then $\nu(h_{\what L}) = \nu(h_L) + 2$.}
\end{corollary}





We now embark on the proof of the bicycle theorem, which will occupy us for the rest of this section.  

%%%%%%%%%%%%%%%%%%%%%%%%%%%
%%%%%%%%%%%%%%%%%%%%%%%%%%%
\part{Proof of the bicycle theorem}
%%%%%%%%%%%%%%%%%%%%%%%%%%%
%%%%%%%%%%%%%%%%%%%%%%%%%%%

Start with a generic link $L$ in $S^3$ where, as usual, $x=x(s)$, $y=y(t)$ and $z=z(u)$ parametrize the three components $X$, $Y$ and $Z$.   Then $Z$ coincides with the binding $K$ of the standard open book, which is the circle subgroup of $S^3$ containing the quaternion $k$.  Looking back at the formula for the characteristic map 
$$
h_L(s,t,u) \ = \ (y_z-x_z)/|y_z-x_z| \quad\text{where}\quad a_z = \pr_{-1}(-a\zbar)\,,
$$
we see that we must understand the $K$-action sending $z\in K$ to the automorphism $a\mapsto a\zbar$ of $S^3$, and the induced $K$-action on $\br^3$ via stereographic projection.  Visualizing how these actions transform the pages of the standard open books will aid us in our subsequent arguments, and we do this next. 



%%%%%%%%%%%%%%%%%%%%%%%%%%%
%%%%%%%%%%%%%%%%%%%%%%%%%%%
\part{The geometry of the $K$-action}
%%%%%%%%%%%%%%%%%%%%%%%%%%%
%%%%%%%%%%%%%%%%%%%%%%%%%%%


For any $z =  \cos\alpha + k\sin\alpha$ in the binding $K$,
right multiplication by $\zbar =  \cos\alpha -k\sin\alpha$ is an isometry of $S^3$ that rotates $K$ by $-\alpha$ radians, likewise rotates the orthogonal great circle $C$ of page centers by $\alpha$ radians, and so advances each hemispherical page $H_\theta$ to the page $H_{\theta+\alpha}$ while simultaneously rotating the page by $-\alpha$ radians about its center $i_\theta$.  Therefore, as $z$ traverses $K$, any given page turns once around $K$, successively occupying the positions vacated by the other pages.  During this time, the page spins once negatively about its center, so that in total it is following a left-handed screw motion along $C$.   

This turning of the hemispherical pages about $K$ in $S^3$ is transferred by stereographic projection to a turning of the half-planar pages about the $k$-axis in $\br^3$, while the {\it spherical} rotations of the pages in $S^3$ become {\it hyperbolic} rotations of the pages in $\br^3$ about their centers.  This is a consequence of the conformality of stereographic projection, which implies that the page identification $H_\theta \leftrightarrow P_\theta$ is conformal.  Once again, the net effect is a left-handed screw motion along $C$.

This description of the $K$-action has the following technical consequence that is critical for our study of the characteristic map of a generic link.

\break

  %%%%%  TWIST LEMMA 5.2  %%%%%
\begin{lemma}\label{lem:twist}
\textbf{\boldmath \textup{(Twist Lemma)}
Let $x$ and $y$ be distinct points in $S^3$ lying in the complement of the binding $K$ of the standard open book.
\begin{indentation}{-.2em}{1.2em}
\begin{itemize}
\item[(a)]  If $x$ and $y$ lie on different pages, then for $z\in K$, the vector $y_z-x_z$ never lies on a page of the corresponding open book in $\br^3$, and in particular never points straight up.
\smallskip  
\item[(b)] If $x$ and $y$ lie on the same page, then as $z$ traverses $K$, the vectors $y_z-x_z$ lie on successive pages in $\br^3$, turning once positively around the binding, and spinning once counterclockwise without backtracking within the pages as they go.  In particular, $y_z-x_z$ points straight up for a unique $z = \tau(x,y) \in K$.    
\end{itemize}
\end{indentation}}
\end{lemma}
%%%%%%%%%%%%%%%%%

\begin{proof}
Since the $K$-action carries pages to pages, $x$ and $y$ will lie on the same page in $S^3$ if and only if $x_z$ and $y_z$ lie on the same page in $\br^3$ for all $z\in K$.  Part (a) of the lemma is now obvious.   Using the meridional projection, and the description of the $K$-action above, part (b) of the lemma translates into the following statement.  For any pair of distinct points $x$ and $y$ in the upper half-plane model of the hyperbolic plane $\bh$, the function 
$$
d(\alpha) \ = \ \arg(\rot_\alpha y-\rot_\alpha x)
$$
is strictly increasing, where $\rot_\alpha:\bh\to\bh$ is hyperbolic rotation about $i$ by $\alpha$ radians.   Since we wish to prove this for all $x$ and $y$, it suffices to show that $d'(0)>0$ (because $d'(\alpha)$ for one choice of $x$ and $y$ is equal to $d'(0)$ for some other choice of $x$ and $y$).    

To prove this, we transfer the problem to the Poincar\'e disk $\bd$ using the conformal map 
$$
f:\bd\ \longrightarrow\ \bh \quad,\quad f(z) \ = \  \frac{1-z}{1+z}\ i
$$
that sends $0$ to $i$.  Hyperbolic rotation of $\bh$ about $i$ by any angle is conjugate by $f$ to euclidean rotation of $\bd$ by the same angle.   Thus we must show that for any pair of distinct points $x$ and $y$ in $\bd$, the function
$$
g(\alpha) \ = \ \arg(f(e^{i\alpha}y)-f(e^{i\alpha}x))
$$
has positive derivative at $\alpha = 0$.   Noting that  $f(a)-f(b) = 2i(b-a)/(1+a)(1+b)$ and using the fact that the argument function converts products into sums and quotients into differences, we find that $g(\alpha)$ differs by a constant from the function 
$$
h(\alpha) \ = \ \alpha - \arg(1+e^{i\alpha}x) - \arg(1+e^{i\alpha}y),
$$  
and so it remains to show that $h'(0)>0$.  But a simple geometric argument using the central angle theorem from elementary plane geometry shows that, for any given $z\in \bd$, the derivative of the function $\arg(1+e^{i\alpha}z)$ at $\alpha = 0$ is strictly less than $1/2$.  Therefore $h'(0) > 1 - 1/2-1/2 \, = \, 0$ as desired.
\end{proof}

It follows from the twist lemma that, for any point $(s,t)$ in the $2$-torus that is isogonal for our generic link $L$ (meaning that  $x=x(s)$ and $y=y(t)$ have the same polar angle, and hence lie on the same page of the open book),  there exists a unique $u = u(s,t)\in S^1$ for which the vector $y_z-x_z$ points straight up, namely the $u$ for which $z = z(u) = \tau(x,y)$.   It follows that the link
$$
\bl \ = \ h_L^{-1}(k) \ = \ \{(s,t,u)\in\tor \st (s,t)\in\cald \text{ and } u=u(s,t)\}\,, \label{pontryaginlink}
$$
is the graph of the function $u(s,t)$ over the collection $\cald$ of isogonal curves in $\torh$, as asserted in the introduction.  

We now use the twist lemma to prove part (a) of the bicycle theorem, asserting that $k$ is a regular value of $h_L$.  This will endow the components of $\bl$ with orientations (and thus vertical winding numbers) and framings, as defined in \secref{pontryagin}.  The proof of the bicycle theorem will then be completed by showing that these agree with the preferred orientations, vertical winding numbers and framings of the components of $\cald$, as defined earlier in this section.

%%%%%%%%%%%%%%%%%%%%%%%%%%%
%%%%%%%%%%%%%%%%%%%%%%%%%%%
\part{Why is $k$ a regular value of $h_L$?}
%%%%%%%%%%%%%%%%%%%%%%%%%%%
%%%%%%%%%%%%%%%%%%%%%%%%%%%

Suppose that $h_L(s,t,u) = k$.  This means that the vector from $x_z $ to $y_z$ points straight up, and so lies in some page $P$ of the standard open book in $\br^3$.
% \foot{\,In particular $P=P_{\theta+\alpha}$ where $\theta$ is the polar angle of $x$ and $y$, and $z = \cos\alpha + k\sin\alpha$.\label{footnote}}  
We must show that at $(s,t,u)$ the vectors $\partial_s h_L$, $\partial_t h_L$ and $\partial_u h_L$ span the tangent space to $S^2$ at $k$.  

We are guided by \figref{regular}, depicting a neighborhood of the vertical vector $\vec$ from $x_z $ to $y_z$ in $\br^3$, lying on the page $P$.   In this figure, $X_z$ and $Y_z$ denote the images of $X$ and $Y$ under the $z$-action, one or both of which must be transverse to $P$ since $L$ is generic.  We arbitrarily depict $X_z$ tangent to $P$ and $Y_z$ transverse to it.  

%%% FIGURE 26: Regular Value %%%
\begin{figure}[h!]
\includegraphics[height=170pt]{FigRegular.pdf}
\put(-114,105){\small$y_z$}
\put(-113,66){\small$x_z$}
\put(-113,83){\small$\vec$}
\put(-21,116){$Y_z$}
\put(-17,84){$X_z$}
\put(-100,120){\tiny$\partial_u h_L$}
\put(-85,92){\tiny$\partial_t h_L$}
\put(-120,29){$P$}
\put(-95,5){\scriptsize TURN}
\put(-56,158){\scriptsize S}
\put(-56.5,150){\scriptsize P}
\put(-55.5,142){\scriptsize I}
\put(-57,134){\scriptsize N}
\caption{$k$ is regular}
\label{fig:regular}
\end{figure}
%%%%%%%%%%%%%%%%%%%%

The hollow arrows in the figure indicate the left-handed screw motion of the $K$-action.  In particular, as the $u$ parameter increases along $K$, this action (the stereographic image of right multiplication by $-\zbar$) turns the pages in the indicated direction while spinning them in a left-handed fashion about their centers. 

As the $t$ parameter along $Y$ increases, the tip of $\vec$ moves to the right along $Y_z$, so that $\partial_t h_L$ also points to the right, transverse to $P$.  For clarity, we draw this partial derivative vector at the tip of $\vec$, even though it is really tangent to $S^2$ at its north pole $k$. 

As the $u$ parameter along $Z$ increases, two things happen which affect $\vec$.  The page containing it turns around the binding, but such an action keeps $\vec$ vertical, so has no infinitesimal effect.  In addition, the page spins around its center, shifting the binding in a downward direction.  Thus by the twist lemma, the vector $\partial_u h_L$ is nonzero  and tangent to the page, pointing toward the binding when placed with its initial point at $y_z$.     

It follows that $\partial_t h_L$ and $\partial_u h_L$ are linearly independent, confirming that $(s,t,u)$ is a regular point of the map $h_L$.  Since $(s,t,u)$ was chosen arbitrarily in $h_L^{-1}(k)$, we see that $k$ is a regular value of $h_L$, proving part (a) of the Bicycle Theorem.

%%%%%%%%%%%%%%%%%%%%%%%%%%%
%%%%%%%%%%%%%%%%%%%%%%%%%%%
\part{Finishing the proof of the Bicycle Theorem}
%%%%%%%%%%%%%%%%%%%%%%%%%%%
%%%%%%%%%%%%%%%%%%%%%%%%%%%

Starting with a three-component link $L$ in generic position in $S^3$, we have shown above 
that the Pontryagin link $\bl$ of its characteristic map $h_L$ is given by
$$
\bl \ = \ h_L^{-1}(k) \ = \ \{(s,t,u)\in\tor \st (s,t)\in\cald \text{ and } u=u(s,t)\}\,. 
$$
It remains to prove that the preferred orientation, vertical winding number $r_i$ and framing $n_i$ of each isogonal curve $\cald_i$, as specified in \figref{local} and just before the statement of the bicycle theorem,  agree with the orientation, vertical winding number and framing of the component $\bl_i$ of the Pontryagin link that lies over $\cald_i$. 

To see that the orientations agree, choose a regular page vector $(x,y)$ parametrized by an isogonal point $(s,t)$ in $\cald_i$, and consider the orientation of $Y$ relative to the page containing $x$ and $y$, as recorded by the sign of $y$.  

\figref{orientation} shows the associated vertical vector $(x_z,y_z)$ in $3$-space for the case when $\sgn(y) = +1$,  causing the displaced link component $Y_z$ to point to the {\it right}.  Arguing as above (see \figref{regular}) we see that  $\partial_t h_L$ also points to the right, while $\partial_u h_L$ points toward the binding as always.  Thus $\partial_t h_L$ and $\partial_u h_L$ form a positive basis for the tangent plane to $S^2$ at $k$.  Since $\partial_s$,  $\partial_t$ and $\partial_u$ form a positive basis for the tangent space to $\tor$ at each of its points, it follows that the link component $\bl_i$ must be oriented in the direction of {\it increasing} $s$ near the associated point $(s,t,u)$ in the $3$-torus, and this agrees with the orientation assigned to $\cald_i$ in \figref{local}.  

%%% FIGURE 27: Orientation %%%
\begin{figure}[h!]
\includegraphics[height=170pt]{FigOrientation.pdf}
\put(-100,100){\small$y_z$}
\put(-100,48){\small$x_z$}
\put(-19,110){$Y_z$}
\put(-17,50){$X_z$}
\put(-81,113){\tiny$\partial_u h_L$}
\put(-66,79){\tiny$\partial_t h_L$}
\caption{Orientation of the Pontryagin link}
\label{fig:orientation}
\end{figure}
%%%%%%%%%%%%%%%%%%%%

Similarly, if $\sgn(y)=-1$ then $\bl_i$  points in the direction of {\it decreasing} $s$ near $(s,t,u)$.  Referring again to \figref{local}, it is evident that this agrees with the preferred orientation on $\cald_i$.  

\begin{comment}
Alternatively, one can show that $\sgn(x)=+1$ or $-1$ according to whether $\bl_i$ is oriented in the direction of increasing or decreasing $t$, which also agrees with the preferred orientation on $\cald_i$.
\end{comment}

Next consider the vertical winding number $r_i$ of $\cald_i$.  By definition $r_i$ is equal to the meridional degree $m_i$ of the associated bicycle $\calp_i$, that is, the number of times that the page vectors spin in the pages of the open book as the bicycle is traversed.  By the twist lemma, this is equal to the number of times that $z$ spins around $K$, or equivalently that $u = u(s,t)$ spins around the circle as $(s,t)$ traverses $\cald_i$, which is the vertical winding number of $\bl_i$. 

Finally consider the framing $n_i$ of $\cald_i$.  By definition $n_i = -\ell_i-m_i$, where $\ell_i$ and $m_i$ are the longitudinal and meridional degrees of the associated bicycle $\calp_i$.
This integer specifies a normal vector field $\vecn$ to $\bl_i$ by adding $n_i$ full twists (positive or negative according to the sign of $n_i$) to the lift $\vecz$ of a normal vector field to $\cald_i$ in $T^2$, the ``zero" or ``blackboard" framing.  

We must show that the vector field $\vecn$ coincides with the Pontryagin framing of $\bl_i$.  In other words, the differential of $h_L$ carries it onto a homotopically trivial loop of nonzero tangent vectors to $S^2$ at $k$.  To see this, first note that $\vecz$ is homotopic to the vertical vector field $\partial_u$, remaining nonzero and transverse to $\bl_i$ during the homotopy.  Now since $\partial_u h_L$ is always tangent to the page containing $x_z$ and $y_z$, as shown above, it is clear that each longitudinal circuit of the bicycle $\calp_i$ causes $dh_L(\vecz)$ to spin once in the same direction about the binding.  By the twist lemma, each meridional spin of the bicycle also causes $dh_L(\vecz)$ to spin once about the binding.  Thus $dh_L(\vecz)$ spins $ \ell_i+m_i$ times as the bicycle is traversed.  Adding $n_i= -\ell_i-m_i$ full twists to the framing negates this spinning, and so $dh_L(\vecn)$ does not spin at all, up to isotopy, as asserted. 

This completes the proof of the bicycle theorem.  



















































